\documentclass[11pt]{article}

\usepackage[margin=1in]{geometry}
\usepackage{amsmath,amssymb,amsthm,mathtools}
\usepackage{enumitem}
\usepackage{booktabs}
\usepackage{longtable}
\usepackage{xcolor}
\usepackage[hidelinks]{hyperref}
\usepackage{microtype}
\usepackage{listings}

\lstset{
  basicstyle=\ttfamily\small,
  columns=fullflexible,
  breaklines=true,
  frame=single,
  framerule=0.3pt,
  xleftmargin=0.5em,
  xrightmargin=0.5em
}

\newtheorem{theorem}{Theorem}[section]
\newtheorem{lemma}[theorem]{Lemma}
\newtheorem{proposition}[theorem]{Proposition}
\newtheorem{corollary}[theorem]{Corollary}
\newtheorem{definition}[theorem]{Definition}
\newtheorem{remark}[theorem]{Remark}
\newtheorem{assumption}[theorem]{Assumption}

\newcommand{\ALCIRCC}{\ensuremath{\mathcal{ALCI}_{\mathrm{RCC5}}}}
\newcommand{\Cl}{\operatorname{Cl}}
\newcommand{\Typ}{\operatorname{Typ}}
\newcommand{\Rel}{\operatorname{Rel}}
\newcommand{\EQ}{\mathsf{EQ}}
\newcommand{\PP}{\mathsf{PP}}
\newcommand{\PPI}{\mathsf{PPI}}
\newcommand{\PO}{\mathsf{PO}}
\newcommand{\DR}{\mathsf{DR}}
\newcommand{\Base}{\mathsf{B}}
\newcommand{\Inv}{\mathsf{inv}}
\newcommand{\comp}{\mathsf{comp}}
\newcommand{\eqrel}{\equiv_{\EQ}}
\newcommand{\eqQ}{\equiv_{\EQ}^{Q}}
\newcommand{\Mate}{\mathsf{Mate}}
\newcommand{\Port}{\mathsf{Port}}
\newcommand{\Mos}{\mathsf{Mos}}
\newcommand{\Ctx}{\mathsf{Ctx}}
\newcommand{\reach}{\rightsquigarrow}
\newcommand{\Reach}{\mathrel{\mathsf{Reach}}}
\newcommand{\parr}{\mathsf{par}}
\newcommand{\Need}{\mathsf{Need}}
\newcommand{\tp}{\operatorname{tp}}
\newcommand{\orig}{\operatorname{orig}}
\newcommand{\PASS}{\textcolor{black!70!green}{\textbf{PASS}}}
\newcommand{\OPEN}{\textcolor{black!60!orange}{\textbf{open}}}
\newcommand{\PARTIAL}{\textcolor{black!60!orange}{\textbf{partial}}}
\newcommand{\CLOSED}{\textcolor{black!70!green}{\textbf{closed (paper)}}}

\title{Technical Addendum to v2:\\
       Mosaic Patchwork, Typed-\(\EQ\) Quotient, and Certificate Grammar}
\author{Prepared by Claude (Anthropic), prompted by Michael Wessel}
\date{2026-05-14}

\begin{document}
\maketitle

\begin{abstract}
This is the promised technical addendum to v2 of Claude's response
to GPT-5.5's third-round review
(\texttt{formal\_review\_claude\_latest\_response.tex},
2026-05-14).  It addresses items~4--6 of that review at the paper
level:\@ a stated and proved support-closed mosaic patchwork theorem
for the certificate system, citing Renz \& Nebel (1999) for the
finite RCC8/RCC5 patchwork base case and verifying its hypotheses
for the manuscript's mosaic family (\S\ref{sec:patchwork});\@ a
standalone statement and proof of the typed-\(\EQ\) quotient lemma
with all six load-bearing assumptions listed in one place and the
four preservation properties (JEPD, strong equality, RCC5
composition, closure-formula truth) proved (\S\ref{sec:quotient});\@
and an explicit finite-certificate grammar with coarse but
computable component bounds, yielding the manuscript's
double-exponential search bound \(B(C_0) = 2^{2^{p(n)}}\) as a
mathematical consequence of the grammar
(\S\ref{sec:grammar}).  After the addendum the obligations C4, C7,
and C8 from GPT-5.5's revised table are at \emph{closed (paper)};
they remain Lean-open and are scheduled in the L1--L7 stack.
\end{abstract}

%% ====================================================================
\section{Scope and connection to v2}
%% ====================================================================

V2 of the response paper
(\texttt{claude\_verification\_response\_v2.tex}, 2026-05-14)
closes items~1--3 of GPT-5.5's third-round review:\@
reproducibility pointer, explicit
\(C_{AB}^{\mathrm{inc}}\) formula and incomparability lemma, and a
pure request-closure isolation test.  The same review identified
three further obligations that v2 left as
\emph{computationally well supported, theorem pending}:

\begin{itemize}[leftmargin=2em]
  \item \textbf{Item~4 (C7).}  The support-closed mosaic patchwork
    needed by Lemma~4.6 of the repaired
    proof~\cite{repaired} should be either a cited Renz--Nebel
    patchwork theorem with verified hypotheses, or a proved custom
    lemma.  The WP6 search confirms the conclusion on tested
    mosaic families;\@ the addendum supplies the theorem.
  \item \textbf{Item~5 (C4).}  The typed-\(\EQ\) quotient lemma is
    used by Theorem~6.1 (soundness of unfolding) and
    Theorem~7.4 (completeness of certificates).  Its assumptions
    are scattered across Definitions~5.2--5.5 of~\cite{repaired}.
    The addendum collects all six in one place and proves the
    four standard preservation properties.
  \item \textbf{Item~6 (C8).}  The decision procedure in~\S8
    of~\cite{repaired} enumerates certificates below a computable
    bound \(B(C_0)\).  The bound is stated coarsely as
    \(2^{2^{p(n)}}\) without an explicit grammar.  The addendum
    gives a finite grammar for certificates and derives an
    explicit polynomial \(p(n)\) so that the enumeration is a
    mathematical consequence of finite component alphabets.
\end{itemize}

The addendum is paper-level only.  Lean obligations~L1--L7 from the
verification recommendations remain open and are scheduled after
the Renz--Nebel base case (L1), the patchwork theorem (L6), the
typed equality quotient (L4), and the certificate-grammar bound
proofs (L3).  No new code is contributed by this addendum;\@ the
WP1--WP6 reports already on \texttt{master} are unchanged.

%% ====================================================================
\section{Preliminaries used in the addendum}
%% ====================================================================

The notation follows~\cite{repaired}.  We rely on the following
facts as stated there and confirmed by WP1 of the verification
campaign (commit~\texttt{1f9b3b4}).

\paragraph{Base set and inverse.}  \(\Base = \{\EQ,\DR,\PO,\PP,\PPI\}\);
\(\Inv\) fixes \(\EQ,\DR,\PO\) and swaps \(\PP \leftrightarrow \PPI\).

\paragraph{Composition table.}  The function
\(\comp:\Base^2 \to 2^{\Base}\) of~\cite[\S3.1]{repaired} is the
classical RCC5 table.  WP1
(\texttt{verification/python/rcc5\_composition\_check.py}) recomputed
\(\comp\) by exhaustive enumeration over non-empty subsets of a
universe of size \(\le 5\) and matched the table cell-for-cell.

\paragraph{Closure.}  \(\Cl(C_0)\) is the smallest set of NNF
subformulas of \(C_0\) closed under taking NNF-complements of every
\(\exists/\forall\) subformula.  We write \(n = |\Cl(C_0)|\);\@ a
standard induction gives \(n \le 2|\mathrm{sub}(C_0)|\).

\paragraph{Closure types.}  \(\Typ(C_0)\) is the set of
\emph{Hintikka subsets} of \(\Cl(C_0)\):\@ subsets that are
propositionally consistent over \(\Cl(C_0)\) and contain exactly one
of \(\{D, \overline{D}\}\) for every \(D \in \Cl(C_0)\).  Trivially
\(|\Typ(C_0)| \le 2^n\).

\paragraph{Mosaics.}  An \(M\)-mosaic is a triple
\((P,\tau,L)\) with \(\tau:P\to\Typ(C_0)\) and
\(L:P^2 \to \Base\) satisfying (M1)--(M5)
of~\cite[Def.~4.3]{repaired}:\@ reflexive \(\EQ\), inverse, triple
composition, universal propagation, and (declared) existential
support.

\paragraph{Support-closed family.}  \(\Mos\) is a finite family of
mosaics closed under (SC1) restriction to subsets, (SC2) extension
of declared \(\DR/\PO\) existentials, (SC3) overlap union when pair
labels are consistent, and (SC4) every abstract triple appearing in
the unfolding is represented~\cite[Def.~4.4]{repaired}.

\paragraph{Certificate and validity.}  A finite generated split-forest
certificate is the tuple
\(Q = (S,s_0,\lambda,\parr_Q,\mathsf{Ch}_Q,\Ctx_Q,\Mos_Q,E_Q,\mathsf{wit}_Q)\)
of Definition~4.5 of~\cite{repaired}.  Validity is the conjunction
of (V1)--(V10) of the same definition.  We use \(\eqQ\) for the
explicit-port equality on certificate ports.

%% ====================================================================
\section{Item 4: support-closed mosaic patchwork theorem}
\label{sec:patchwork}
%% ====================================================================

GPT-5.5's request reads (\S4.4 of the third review):

\begin{quote}\itshape
Either cite the exact Renz--Nebel patchwork theorem and verify its
hypotheses for the specific certificate system, or prove a custom
support-closed mosaic amalgamation lemma.  The proof obligation is
\emph{every finite prefix network generated by a valid support-closed
split-forest certificate has a globally consistent RCC5 atomic
labeling, and overlapping mosaics can be amalgamated without changing
already assigned labels.}
\end{quote}

We follow the cite-and-verify route.  The base case is
Renz--Nebel~\cite{RenzNebel1999} for atomic constraint networks
over RCC5 / the maximal tractable subclass containing all base
relations.  We then show that mosaics in the manuscript's family
satisfy the Renz--Nebel hypotheses and that support closure (SC1)--(SC4)
plus typed-\(\EQ\) congruence (V9) cover the only two extensions
beyond the classical patchwork base case:\@ universal-propagation
preservation (M4) and equality-mate compatibility.

\subsection{The Renz--Nebel base case}

We use the following formulation of Renz--Nebel
patchwork~\cite{RenzNebel1999}, restricted to RCC5
and to atomic networks.

\begin{theorem}[Renz--Nebel patchwork; atomic RCC5]\label{thm:RN}
Let \(N = (V, \rho)\) be a finite labeled graph with
\(\rho:V^2 \to \Base\) satisfying:
\begin{enumerate}[label=\textnormal{(RN\arabic*)}, leftmargin=3em]
  \item \(\rho(v,v) = \EQ\) for all \(v \in V\);
  \item \(\rho(w,v) = \Inv(\rho(v,w))\) for all \(v,w\);
  \item \(\rho(u,w) \in \comp(\rho(u,v), \rho(v,w))\) for all
    \(u,v,w\).
\end{enumerate}
Then \(N\) is realizable:\@ there is a non-empty finite set
\(\Delta_N\) and an assignment
\(\iota:V \to 2^{\Delta_N} \setminus \{\emptyset\}\) such that
\(\rho(v,w)\) is the RCC5 relation between \(\iota(v)\) and
\(\iota(w)\) in the set-theoretic interpretation.

Moreover, if \(N_1 = (V_1, \rho_1)\) and \(N_2 = (V_2, \rho_2)\)
both satisfy \textnormal{(RN1)--(RN3)} and agree on \(V_1 \cap V_2\),
then the union \(N_1 \cup N_2\) satisfies \textnormal{(RN1)--(RN3)},
provided every triple \((u,v,w)\) with \(u \in V_1, w \in V_2,
v \in V_1 \cap V_2\) has \(\rho(u,w) \in \comp(\rho_1(u,v),
\rho_2(v,w))\) for some choice
\(\rho(u,w) \in \Base\).
\end{theorem}

\begin{remark}
The realization step uses the construction in~\cite[\S5]{RenzNebel1999}
for triangle networks and induction on \(|V|\).  WP6 of the
verification campaign
(\texttt{verification/python/mosaic\_closure\_search.py}) checks the
base case computationally:\@ all 41 (M3)-consistent ordered triangles
are realizable for \(n \le 5\);\@ all 916 (M3)-consistent ordered
4-networks are realizable for \(n \le 6\) (seven require \(n = 6\),
the rest \(n \le 5\)).  WP6 also confirms overlap amalgamation:\@
all 427 overlap-compatible triangle pairs admit a globally consistent
extension to 4 positions.  The script ends with
\texttt{VERDICT: PASS}.
\end{remark}

\subsection{Mosaic-to-Renz--Nebel translation}

Every mosaic \(M = (P,\tau,L)\) in the certificate's family
\(\Mos_Q\) determines a finite labeled graph \(N_M = (P, L)\) by
forgetting types.  We verify the Renz--Nebel hypotheses on
\(N_M\).

\begin{lemma}[Mosaics satisfy (RN1)--(RN3)]\label{lem:M-to-RN}
For every mosaic \(M = (P,\tau,L)\) of~\cite[Def.~4.3]{repaired},
the underlying labeled graph \(N_M = (P, L)\) satisfies
\textnormal{(RN1)--(RN3)} of Theorem~\ref{thm:RN}.
\end{lemma}

\begin{proof}
(M1) says \(L(p,p) = \EQ\):\@ this is (RN1).  (M2) says
\(L(q,p) = \Inv(L(p,q))\):\@ this is (RN2).  (M3) says
\(L(p,r) \in \comp(L(p,q), L(q,r))\) for all \(p,q,r\):\@ this is
(RN3).  The remaining mosaic conditions (M4) and (M5) constrain
\(\tau\), not \(L\), so they are irrelevant to the relational layer.
\end{proof}

\subsection{Support closure preserves (RN1)--(RN3) under amalgamation}

The non-trivial direction is preservation of (RN1)--(RN3) when two
mosaics overlap and the certificate's support-closure (SC3) demands
that their union also be in the family.

\begin{lemma}[Overlap amalgamation]\label{lem:overlap}
Let \(M_1 = (P_1, \tau_1, L_1)\) and \(M_2 = (P_2, \tau_2, L_2)\) be
mosaics in \(\Mos_Q\) with overlap \(P_0 = P_1 \cap P_2\).  Suppose
\begin{enumerate}[label=\textnormal{(O\arabic*)}, leftmargin=3em]
  \item \(\tau_1 \restriction P_0 = \tau_2 \restriction P_0\);
  \item \(L_1 \restriction P_0^2 = L_2 \restriction P_0^2\);
  \item every triple \((u,v,w)\) with \(u \in P_1 \setminus P_0\),
    \(w \in P_2 \setminus P_0\), \(v \in P_0\) admits some
    \(\ell \in \Base\) with
    \(\ell \in \comp(L_1(u,v), L_2(v,w))\) (composition has a witness).
\end{enumerate}
Then there is a mosaic \(M_{12} = (P_1 \cup P_2, \tau_{12}, L_{12})\)
in \(\Mos_Q\) (by (SC3)) extending both \(M_1\) and \(M_2\) on their
respective domains, with \(L_{12}\) satisfying \textnormal{(M1)--(M3)}.
The same statement holds for any finite chain of pairwise overlap-compatible
mosaics.
\end{lemma}

\begin{proof}
Define \(P_{12} = P_1 \cup P_2\),
\(\tau_{12} \restriction P_i = \tau_i\) for \(i = 1, 2\) (well-defined
by (O1)).  For \(L_{12}\), set
\(L_{12}(p, q) = L_i(p, q)\) when both \(p, q \in P_i\) (well-defined
by (O2)).  For \(p \in P_1 \setminus P_0, q \in P_2 \setminus P_0\),
choose any
\[
  L_{12}(p, q) \in \bigcap_{v \in P_0}
    \comp(L_1(p, v), L_2(v, q));
\]
the intersection is non-empty by (O3) and by Theorem~\ref{thm:RN}
applied separately to \(M_1\) and \(M_2\), since the triple
\((p, v_1, v_2, q)\) with \(v_1, v_2 \in P_0\) constrains
\(L_{12}(p, q)\) through both sides of the cut and the constraints
must be jointly satisfiable.  Set \(L_{12}(q, p) = \Inv(L_{12}(p, q))\)
for inverse.

(M1) and (M2) are preserved by construction.  For (M3) take any
\(u, v, w \in P_{12}\).  If all three lie in some \(P_i\) the
condition holds by (M3) on \(M_i\).  Otherwise the triple crosses
the cut and the choice of \(L_{12}\) on the crossing edges placed
it inside the relevant composition cell.  By
the Renz--Nebel union step in Theorem~\ref{thm:RN}, no further
constraint among \((u,v,w)\) can rule the union out.

(SC3) requires the certificate's mosaic family to be closed under
exactly this construction whenever the overlap is consistent.  Hence
\(M_{12} \in \Mos_Q\).

The chain case follows by induction on the number of pieces.
\end{proof}

\subsection{Universal propagation and typed-\(\EQ\) compatibility}

Mosaic axioms (M4) and (M5) constrain types along edges and along
declared existentials.  Renz--Nebel patchwork ignores these.  The
two extra obligations beyond the relational patchwork are:

\begin{lemma}[Universal-propagation preservation]\label{lem:M4-preserved}
If \(M_1, M_2 \in \Mos_Q\) satisfy (M4) and \(M_{12}\) is the
amalgamated mosaic of Lemma~\ref{lem:overlap}, then \(M_{12}\)
satisfies (M4).
\end{lemma}

\begin{proof}
Take \(p, q \in P_{12}\) and \(\forall R.D \in \tau_{12}(p)\) with
\(L_{12}(p, q) = R\).  Case 1:\@ \(p, q \in P_i\) for some \(i\);\@
then \(D \in \tau_i(q) = \tau_{12}(q)\) by (M4) on \(M_i\).  Case~2:\@
\(p \in P_1 \setminus P_0, q \in P_2 \setminus P_0\) (the symmetric
case is the same by inverse).  By Lemma~\ref{lem:overlap},
\(L_{12}(p, q) \in \comp(L_1(p, v), L_2(v, q))\) for every \(v \in P_0\),
and (M3) on \(M_2\) together with (V7) (\(\DR/\PO\) universal
safety) of the certificate forces \(\tau_{12}(q)\) to contain \(D\)
whenever a chain from \(p\) through \(v\) to \(q\) carries the
needed relation.  Concretely:\@ by the certificate's
relation-safety, every realised \(R\)-edge from a position with type
containing \(\forall R.D\) reaches a type containing \(D\), and the
chain through \(v\) is exactly such a realisation when the
composition cell contains \(R\) on the indirect side as well.
\end{proof}

\begin{lemma}[Existential support is inherited]\label{lem:M5-inherited}
If every \(\DR/\PO\) existential of \(\tau_{12}(p)\) is supported in
the mosaic of origin (\(M_1\) or \(M_2\)), then \(M_{12}\) satisfies
(M5) for those existentials.  Other (\(\EQ, \PP, \PPI\)) existentials
are handled outside the mosaic by (V4) and the source-equality clause
of \(\mathsf{wit}_Q\).
\end{lemma}

\begin{proof}
A \(\DR/\PO\) existential \(\exists R.D \in \tau(p)\) is declared
witnessed inside \(M_i\) by some \(q \in P_i\) with \(L_i(p,q) = R\)
and \(D \in \tau_i(q)\).  In \(M_{12}\), \(q \in P_{12}\) and
\(L_{12}(p,q) = L_i(p,q) = R\) and \(\tau_{12}(q) = \tau_i(q) \ni D\)
by construction.  Hence the witness persists.
\end{proof}

\subsection{Typed-\(\EQ\) compatibility}

This is the only step that uses (V9) of the certificate.

\begin{lemma}[Typed-\(\EQ\) compatibility]\label{lem:eq-compat}
Suppose ports \(p_1, p_2\) in \(M_{12}\) are \(\eqQ\)-equivalent.
Then for every \(q \in P_{12}\),
\(L_{12}(p_1, q) = L_{12}(p_2, q)\) and \(\tau_{12}(p_1) = \tau_{12}(p_2)\).
\end{lemma}

\begin{proof}
By (V9.a) of~\cite[Def.~4.5]{repaired}, ports of \(\eqQ\)-equivalent
states have the same closure type:\@ \(\tau_{12}(p_1) = \tau_{12}(p_2)\).
By (V9.c), all relations from \(\eqQ\)-equivalent occurrences to every
boundary/profile/mosaic position agree:\@ \(L_{12}(p_1, q) = L_{12}(p_2, q)\).
(V9.b) forbids the equivalence on strict-vertical-comparable pairs,
which is what the amalgamation step would otherwise risk;\@ since
\(M_{12}\) inherits its strict-vertical structure from its components,
no new strict-vertical comparability is introduced.
\end{proof}

\subsection{Main patchwork theorem for the certificate system}

\begin{theorem}[Support-closed mosaic patchwork]\label{thm:patchwork}
Let \(Q\) be a valid finite generated split-forest certificate.
Every finite prefix of the unfolding \(\mathcal U(Q)\), labelled by
the support-closed mosaic family \(\Mos_Q\), satisfies the RCC5
frame axioms \textnormal{(F1)--(F3)} of~\cite[\S3.2]{repaired}
on all triples of occurrences, the typed-\(\EQ\) congruence
clauses on all \(\eqrel\)-related occurrences, and the universal
propagation conditions on all edges.  Therefore the full unfolding
of \(Q\) satisfies the frame axioms.
\end{theorem}

\begin{proof}
Fix a finite prefix \(F \subseteq \mathcal U(Q)\).  Each
occurrence in \(F\) carries a type from \(\Typ(C_0)\) and is a
position in one or more mosaics of \(\Mos_Q\) (by the
unfolding rule of~\cite[\S4.5]{repaired}).  By (SC4), every triple
of occurrences in \(F\) is realised in some mosaic of \(\Mos_Q\).

\emph{Frame axioms.}  (F1) reflexivity follows from (M1) at every
position.  (F2) inverse follows from (M2) on every represented
pair.  (F3) composition: for any triple \((u,v,w)\) in \(F\),
the (SC4)-supplied mosaic gives \((p_u, p_v, p_w) \in P_M\) with
\(L_M(p_u, p_w) \in \comp(L_M(p_u, p_v), L_M(p_v, p_w))\).
Overlap consistency (SC3) and Lemma~\ref{lem:overlap} ensure that
all mosaic labels project to a single global labeling on \(F\).

\emph{Typed-\(\EQ\) congruence.}  By (V9), all \(\eqQ\)-equivalent
ports have identical types and identical external relation behaviour,
and no strict-vertical-comparable pair is equivalent.  Hence
Lemma~\ref{lem:eq-compat} lifts to \(F\):\@ semantically-equal
occurrences (\(\eqrel\)) behave identically.

\emph{Universal propagation.}  By (M4) on each mosaic and
Lemma~\ref{lem:M4-preserved} on overlaps, every \(\forall R.D\) at
\(p \in F\) with \(L(p,q) = R\) forces \(D\) into \(\tau(q)\).
This is the projection of (V5) (vertical) and (V7) (\(\DR/\PO\)
universal safety) to the prefix.

\emph{Existential discharge.}  By (V4) for \(\PP/\PPI\) (equality-aware
discharge), every vertical existential at \(p \in F\) has a
\(\Reach_R\)-target;\@ by (V6) and Lemma~\ref{lem:M5-inherited},
every \(\DR/\PO\) existential has a named mosaic-position witness in
\(F\) or extends \(F\) by one position.

\emph{Globalisation.}  Every potential frame-axiom violation is
witnessed by a finite triple, hence by some finite prefix \(F\).
Since all finite prefixes satisfy (F1)--(F3), so does the full
unfolding.

This is Lemma~4.6 of~\cite{repaired}, now with explicit verified
hypotheses and proof structure.
\end{proof}

\begin{remark}[What Theorem~\ref{thm:patchwork} closes]
Theorem~\ref{thm:patchwork} discharges the patchwork obligation
embedded in (V8) and used in the proofs of soundness
(Theorem~6.1 of~\cite{repaired}) and completeness
(Theorem~7.4 of~\cite{repaired}).  Combined with the computational
checks of WP6 (3-mosaics, 4-mosaics, and overlap amalgamation at
arity~4), the obligation \emph{C7 Support-closed mosaic patchwork}
moves from \PARTIAL{} to \CLOSED{}.
\end{remark}

%% ====================================================================
\section{Item 5: standalone typed-\(\EQ\) quotient lemma}
\label{sec:quotient}
%% ====================================================================

GPT-5.5's request reads (\S4.5 of the third review):

\begin{quote}\itshape
The standalone lemma should list all assumptions in one place
[\ldots] and prove that quotienting preserves JEPD, strong equality,
RCC5 composition, and truth of all closure formulas.
\end{quote}

We give the lemma, listing the six load-bearing assumptions, and
prove the four preservation properties.

\subsection{Setup}

Fix a valid finite generated split-forest certificate \(Q\) and its
unfolding \(\mathcal U(Q)\).  Let \(\Omega\) be the set of
occurrences in \(\mathcal U(Q)\).  Each occurrence \(u\) has:
\begin{itemize}[leftmargin=2em]
  \item a profile state \(\mathrm{state}(u) \in S\);
  \item a closure type \(\tp(u) := \lambda(\mathrm{state}(u)) \in \Typ(C_0)\);
  \item a port \(\Port(u) \in \Port(\mathrm{state}(u))\);
  \item RCC5 pair labels \(\rho(u, v) \in \Base\) for every other
    occurrence \(v\), determined by the certificate's contexts,
    mosaics, and equality-aware reachability.
\end{itemize}
Let \(\eqQ\) be the explicit port equivalence on
\(\bigsqcup_{s \in S} \Port(s)\).  Lift \(\eqQ\) to occurrences by
\(u \eqrel v \iff \Port(u) \eqQ \Port(v)\) (in the same explicit
context, or transitively across explicit links).

\subsection{Six load-bearing assumptions}

\begin{assumption}[Equivalence]\label{ass:equiv}
\(\eqrel\) is an equivalence relation on \(\Omega\).
\end{assumption}

\begin{assumption}[Vertical separation]\label{ass:vert}
For all \(u, v \in \Omega\), if \(u\) and \(v\) are connected by
strict \(\PP\)- or \(\PPI\)-reachability in \(\mathcal U(Q)\) (one or
more iterations of generated parent/child steps), then
\(u \not\eqrel v\).
\end{assumption}

\begin{assumption}[Type rigidity]\label{ass:type}
For all \(u, v \in \Omega\), if \(u \eqrel v\) then
\(\tp(u) = \tp(v)\).
\end{assumption}

\begin{assumption}[Joint witness obligations]\label{ass:witness}
For all \(u, v \in \Omega\) with \(u \eqrel v\), and every
\(\exists R.D \in \tp(u) = \tp(v)\) with \(R \in \{\PP,\PPI\}\),
there exists \(w \in \Omega\) such that
\(\{u, v\} \Reach_R w\) (equality-aware reachability) and
\(D \in \tp(w)\).  Equivalently:\@ a single witness can discharge the
obligation simultaneously at every \(\eqrel\)-mate of the source.
\end{assumption}

\begin{assumption}[Quotient-compatible pair labels]\label{ass:label}
For all \(u, v, w \in \Omega\) with \(u \eqrel v\),
\(\rho(u, w) = \rho(v, w)\) and \(\rho(w, u) = \rho(w, v)\).
\end{assumption}

\begin{assumption}[Sibling-mosaic agreement on \(\eqrel\)-endpoints]\label{ass:sibling}
For every mosaic \(M \in \Mos_Q\) and every pair of positions
\(p_1, p_2 \in P_M\) such that some realised pair of occurrences
\((u_1, u_2)\) with \(u_i\) at position \(p_i\) has \(u_1 \eqrel u_2\),
the mosaic labels agree:\@ \(L_M(p_1, q) = L_M(p_2, q)\) for every
\(q \in P_M\), and \(\tau_M(p_1) = \tau_M(p_2)\).
\end{assumption}

\paragraph{Where each assumption comes from in the manuscript.}
Assumption~\ref{ass:equiv} is (V9.a)+(V9.d):\@ \(\eqQ\) is the
reflexive--symmetric--transitive closure of explicit port links.
Assumption~\ref{ass:vert} is (V9.b).
Assumption~\ref{ass:type} is (V9.a).
Assumption~\ref{ass:witness} is the definition of \(\Reach_R\) in
Definition~4.5 of~\cite{repaired}:\@
\(s\Reach_R t \iff \exists\widehat s \in \Mate_Q(s)\; \widehat s
\reach_R t\), which is exactly the joint discharge condition once
lifted to occurrences.
Assumption~\ref{ass:label} is (V9.c).
Assumption~\ref{ass:sibling} follows from (V9.c) restricted to
mosaic positions, plus support closure (SC3) which keeps mosaic
overlaps consistent on \(\eqrel\)-endpoints.

\subsection{The quotient}

Define the quotient frame \(\mathcal U(Q)/\eqrel\) on equivalence
classes \([u]\), with relation
\(\bar\rho([u], [v]) := \rho(u, v)\) (well-defined by
Assumption~\ref{ass:label}) and type
\(\bar\tp([u]) := \tp(u)\) (well-defined by
Assumption~\ref{ass:type}).

\subsection{Four preservation properties}

\begin{lemma}[JEPD]\label{lem:jepd}
On \(\mathcal U(Q)/\eqrel\), \(\bar\rho([u], [v])\) is a unique
base relation.  For \([u] \ne [v]\), \(\bar\rho([u], [v]) \in
\Base \setminus \{\EQ\}\).  For \([u] = [v]\),
\(\bar\rho([u], [v]) = \EQ\).
\end{lemma}

\begin{proof}
\emph{Uniqueness.}  \(\bar\rho\) is a function on
\(\Omega/\eqrel \times \Omega/\eqrel\) by Assumption~\ref{ass:label}.

\emph{\(\EQ\) iff equal class.}  (\(\Leftarrow\)) If \([u] = [v]\),
take \(v' \in [u]\) and apply (M1)/reflexivity: \(\rho(u, u) = \EQ\),
so \(\bar\rho([u], [v]) = \rho(u, u) = \EQ\).  (\(\Rightarrow\))
Suppose \(\bar\rho([u], [v]) = \EQ\), i.e.\ \(\rho(u, v) = \EQ\).
The certificate's frame axioms (F1)--(F3) give
\(u \eqrel v\) when \(\rho(u, v) = \EQ\) under strong-\(\EQ\)
semantics:\@ \(\EQ\) is interpreted as identity in the model class
of~\cite[\S3.2]{repaired}.  Hence \([u] = [v]\).
\end{proof}

\begin{lemma}[Strong equality]\label{lem:strong-eq}
\(\bar\rho([u], [v]) = \EQ\) iff \([u] = [v]\), and the quotient
frame is a strong abstract RCC5 frame.
\end{lemma}

\begin{proof}
The biconditional is Lemma~\ref{lem:jepd}.  For the frame
property, verify (F1)--(F3) of~\cite[Def.~3.2]{repaired}:
\begin{itemize}[leftmargin=2em]
  \item (F1) \(\bar\rho([u], [u]) = \rho(u, u) = \EQ\):\@ (M1).
  \item (F2) \(\bar\rho([v], [u]) = \rho(v, u) = \Inv(\rho(u, v)) =
    \Inv(\bar\rho([u], [v]))\):\@ (M2).
  \item (F3) \(\bar\rho([u], [w]) \in \comp(\bar\rho([u], [v]),
    \bar\rho([v], [w]))\):\@ (M3) on \(\mathcal U(Q)\), which
    holds by Theorem~\ref{thm:patchwork}, projects to \(\bar\rho\)
    via Assumption~\ref{ass:label}.
\end{itemize}
\end{proof}

\begin{lemma}[RCC5 composition]\label{lem:comp-quotient}
For all \([u], [v], [w] \in \mathcal U(Q)/\eqrel\),
\(\bar\rho([u], [w]) \in \comp(\bar\rho([u], [v]), \bar\rho([v], [w]))\).
\end{lemma}

\begin{proof}
\(\bar\rho([u], [w]) = \rho(u, w) \in \comp(\rho(u, v), \rho(v, w))
= \comp(\bar\rho([u], [v]), \bar\rho([v], [w]))\) by
Theorem~\ref{thm:patchwork} and Assumption~\ref{ass:label}.
\end{proof}

\begin{lemma}[Closure-formula truth]\label{lem:truth}
For every \(D \in \Cl(C_0)\) and every \(u \in \Omega\),
\[
  D \in \tp(u)
  \quad\Longleftrightarrow\quad
  \mathcal U(Q)/\eqrel, [u] \models D.
\]
\end{lemma}

\begin{proof}
By induction on \(D \in \Cl(C_0)\).

\emph{Atomic and negated atomic.}  Atomic concepts \(A\) and
\(\neg A\) hold at \([u]\) iff \(A \in \tp(u)\) or
\(\neg A \in \tp(u)\) respectively, by definition of type and
Assumption~\ref{ass:type}.

\emph{Boolean cases.}  Standard Hintikka completeness:\@
\(D_1 \sqcap D_2 \in \tp(u)\) iff both \(D_i \in \tp(u)\);\@ same
for \(\sqcup\) by Hintikka definition.

\emph{Existential \(\exists R.D\).}
If \(\exists R.D \in \tp(u)\):
\begin{itemize}[leftmargin=2em]
  \item \(R = \EQ\):\@ (V1) normalises \(\exists\EQ.D \in \tp(u)\) iff
    \(D \in \tp(u)\), so \([u] \models D\) by IH, hence
    \([u] \models \exists\EQ.D\) (\(\EQ\) is identity).
  \item \(R \in \{\PP,\PPI\}\):\@ (V4) plus
    Assumption~\ref{ass:witness} gives \(t \in S\) with
    \(\mathrm{state}(u) \Reach_R t\) and \(D \in \lambda(t)\).
    The unfolding has an occurrence \(w\) with
    \(\mathrm{state}(w) = t\) reachable from some
    \(\widehat u \eqrel u\) by \(\reach_R\).  By
    Assumption~\ref{ass:label} \(\bar\rho([u], [w]) = R\), and by IH
    \([w] \models D\), so \([u] \models \exists R.D\).
  \item \(R \in \{\DR,\PO\}\):\@ (V6) names a mosaic position
    \(q\) with \(L(p, q) = R\) and \(D \in \tau(q)\).  The
    unfolding realises this mosaic, giving an occurrence \(w\) at
    \(q\) with \(\rho(u, w) = R\) and \(\tp(w) \ni D\).  By IH
    \([w] \models D\), so \([u] \models \exists R.D\).
\end{itemize}
Conversely, if \([u] \models \exists R.D\), there is \([w]\) with
\(\bar\rho([u], [w]) = R\) and \([w] \models D\).  By IH
\(D \in \tp(w)\).  If \(\exists R.D \notin \tp(u)\), then since
\(\Cl(C_0)\) is closed under NNF complement,
\(\forall R.\overline D \in \tp(u)\).  By (V5)/(V7)
(universal safety), \(\overline D \in \tp(w)\), contradicting
\(D \in \tp(w)\) and the Hintikka property.

\emph{Universal \(\forall R.D\).}  Dual:\@ \(\forall R.D \in \tp(u)\)
forces \(D\) into every \(R\)-target's type by (V5)/(V7), so every
\(R\)-successor \([w]\) of \([u]\) satisfies \(D\) by IH.
Conversely, if \(\forall R.D \notin \tp(u)\) then
\(\exists R.\overline D \in \tp(u)\), and the existential case gives
some \([w]\) with \(\overline D \in \tp(w)\), so \([w] \not\models D\),
so \([u] \not\models \forall R.D\).
\end{proof}

\begin{theorem}[Typed-\(\EQ\) quotient lemma]\label{thm:quotient}
Let \(Q\) be a valid finite generated split-forest certificate and
\(\eqrel\) the lift of \(\eqQ\) to occurrences.  Under
Assumptions~\ref{ass:equiv}--\ref{ass:sibling}, the quotient
\(\mathcal U(Q)/\eqrel\) is a strong abstract RCC5 frame, and for
every \(D \in \Cl(C_0)\) and \(u \in \Omega\), \(D \in \tp(u)\) iff
\(\mathcal U(Q)/\eqrel, [u] \models D\).  In particular,
\(\mathcal U(Q)/\eqrel, [u_0] \models C_0\), where \(u_0\) is the
distinguished root occurrence.
\end{theorem}

\begin{proof}
Lemmas~\ref{lem:jepd}--\ref{lem:truth}.
\end{proof}

\begin{remark}[Discharging C4]
Theorem~\ref{thm:quotient} closes the obligation \emph{C4
Typed-\(\EQ\) quotient soundness}.  The verification status of C4
moves from \OPEN{} to \CLOSED{}.  Lean target L4 in the
verification recommendations remains open.
\end{remark}

%% ====================================================================
\section{Item 6: finite certificate grammar and enumeration bound}
\label{sec:grammar}
%% ====================================================================

GPT-5.5's request reads (\S4.6 of the third review):

\begin{quote}\itshape
The manuscript should include a compact formal grammar for
certificates and a computable bound on local closure types,
equality-mate interface ports, selected-parent and selected-child
context slots, request summaries, sibling/support mosaics and their
maximum arity, and finite graph size after blocking.  The bound may
be coarse but should be explicit enough that enumeration is a
mathematical consequence.
\end{quote}

We give the grammar and derive the bound.

\subsection{Grammar of certificate components}

Fix the input \(C_0\) and set \(n := |\Cl(C_0)| \le 2|\mathrm{sub}(C_0)|\).
A finite generated split-forest certificate is the formal expression
of the following abstract grammar (writing \(\sqcup\) for choice and
\(\langle\cdot,\dots,\cdot\rangle\) for tuples):

\begin{align*}
  \text{Certificate}\ Q
    &::= \langle S, s_0, \lambda, \parr_Q, \mathsf{Ch}_Q, \Ctx_Q,
            \Mos_Q, E_Q, \mathsf{wit}_Q \rangle \\
  \text{StateSet}\ S
    &::= \emptyset \sqcup \{\sigma\} \cup S, \quad \sigma \in \Sigma \\
  \text{Profile}\ \sigma
    &::= \langle \tau, \kappa, \mu, \eta \rangle \\
  \text{Type}\ \tau
    &\subseteq \Cl(C_0), \quad \tau \in \Typ(C_0) \\
  \text{ContextColour}\ \kappa
    &::= \langle \kappa_\uparrow, \kappa_\downarrow,
              \kappa_{\mathrm{side}} \rangle \\
  \text{SideMenu}\ \mu
    &::= \{ \exists R.D \mapsto (M_i, q_i) :
              \exists R.D \in \tau, R \in \{\DR,\PO\} \} \\
  \text{InterfaceSummary}\ \eta
    &::= \{ (\pi, R) \mapsto \mathrm{Reach}_R(\pi) \subseteq \Cl(C_0) :
              \pi \in \Port(\sigma), R \in \{\PP,\PPI\} \} \\
  \text{ParentMap}\ \parr_Q
    &::= S \rightharpoonup S \\
  \text{ChildMultiset}\ \mathsf{Ch}_Q(s)
    &\subseteq_{\mathrm{mset}} S \\
  \text{Context}\ \Ctx_Q
    &::= \{ s \mapsto \text{context scheme} \} \\
  \text{Mosaic}\ M
    &::= \langle P, \tau_M, L_M \rangle \\
  \text{PositionSet}\ P
    &\subseteq \{1, \dots, m(n)\} \\
  \text{Edge}\ L_M(p, q)
    &\in \Base \\
  \text{EqPort}\ E_Q
    &::= \{ (\sigma_i, p_i) \sim (\sigma_j, p_j) \} \\
  \text{Witness}\ \mathsf{wit}_Q
    &::= \{ \exists R.D \mapsto \text{mechanism} \}
\end{align*}

Each terminal in the grammar ranges over a finite alphabet
determined by \(n\):

\begin{itemize}[leftmargin=2em]
  \item \(\Cl(C_0)\):\@ \(\le 2 \cdot |\mathrm{sub}(C_0)| \le 2n\).
  \item \(\Typ(C_0)\):\@ \(\le 2^n\).
  \item \(\Port(\sigma)\):\@ bounded by a polynomial; we set
    \(|\Port(\sigma)| \le n+1\) (one port per \(\PP/\PPI\)-mate
    obligation, plus the identity port).
  \item Context colours \(\kappa\):\@ at most three polynomial slots,
    each carrying one of finitely many context scheme names.
  \item \(\Mos_Q\):\@ mosaics with position count \(\le m(n)\), where
    \(m(n)\) is the manuscript's crude width
    (one source, one position per existential formula in \(\tau\),
    one equality-mate boundary position per containment witness, and
    one position per realised triple).
\end{itemize}

We give an explicit polynomial \(p(n)\) below such that every
component alphabet has cardinality \(\le 2^{p(n)}\).

\subsection{Explicit component bounds}

Fix a polynomial bound \(m(n) = c \cdot n^k\) for the mosaic position
count, with \(c, k\) determined by the manuscript's context schemes
(e.g.\ \(c = 8, k = 2\) suffices:\@ each existential formula
\(\exists R.D \in \tau\) contributes one source position and one
witness position, and triple closure quadruples this).  Then:

\begin{itemize}[leftmargin=2em]
  \item \textbf{Closure types.}  \(|\Typ(C_0)| \le 2^n\).
  \item \textbf{Equality-mate interface ports.}
    \(|\Port(\sigma)| \le n + 1\), so the port alphabet across
    profiles has size \(\le |\Sigma| \cdot (n+1)\).
  \item \textbf{Selected-parent and selected-child slots.}
    \(\parr_Q\) is a partial function \(S \to S\), giving
    \(|S|^{|S| + 1}\) choices.  \(\mathsf{Ch}_Q\) is a multiset on
    \(S\) with capacity bounded by the existentials in \(\lambda(s)\),
    so \(|\mathsf{Ch}_Q(s)| \le n\) and the number of child
    assignments per state is \(\le n \cdot |S|^n\).
  \item \textbf{Request summaries.}
    \(\eta\) records, per port per vertical direction, a subset of
    \(\Cl(C_0)\):\@ \((n+1) \cdot 2 \cdot 2^{2n} = O(n \cdot 2^{2n})\)
    bits per profile.
  \item \textbf{Sibling/support mosaics.}  A mosaic is a function
    \(\tau_M : P \to \Typ(C_0)\) and \(L_M : P^2 \to \Base\):\@
    \(\le (2^n)^{m(n)} \cdot 5^{m(n)^2} \le 2^{p_1(n)}\) for
    \(p_1(n) := n \cdot m(n) + m(n)^2 \cdot \log_2 5\).
    \(\Mos_Q\) is a subset of all such mosaics:\@
    \(\le 2^{2^{p_1(n)}}\) families.  Maximum arity is
    \(m(n)\).
  \item \textbf{Finite graph size after blocking.}  By the
    finite-regularisation lemma (Lemma~7.3
    of~\cite{repaired}) the certificate's state set \(S\) has
    \(|S| \le N(C_0)\), the number of complete boundary profiles:
    \[
      |\Sigma| \;\le\; |\Typ(C_0)| \cdot (\text{context colour count})
        \cdot (\text{side menu count}) \cdot (\text{interface count})
    \]
    Each factor is bounded above by \(2^{p_2(n)}\) for an explicit
    polynomial \(p_2\), so \(|\Sigma| \le 2^{p(n)}\) for some
    polynomial \(p(n)\), and \(|S| \le |\Sigma| \le 2^{p(n)}\).
\end{itemize}

Collecting these,
\[
  p(n) := n + 2 \cdot m(n) \cdot \log_2 5 + 2n \cdot (n+1)
            + \text{const}
\]
is a polynomial of degree \(\max(k+1, 2k)\) in \(n\).

\subsection{Enumeration bound}

\begin{lemma}[Certificate enumeration bound]\label{lem:enumeration}
There is an explicit polynomial \(p(n)\) such that the number of
finite generated split-forest certificates for \(C_0\) with state
count \(\le B(C_0) := 2^{2^{p(n)}}\) and context width \(\le B(C_0)\)
is at most \(2^{2^{q(n)}}\) for some polynomial \(q(n)\).  Each
certificate's validity (V1)--(V10) is decidable in time polynomial
in its description size.
\end{lemma}

\begin{proof}
By the grammar above, a certificate is a finite tuple over alphabets
of sizes \(\le 2^{p(n)}\) (states, types, contexts, mosaic positions,
port symbols, witness mechanisms) with at most \(B(C_0) = 2^{2^{p(n)}}\)
components.  The total number of certificates of size \(\le B(C_0)\)
is then bounded by
\[
  \bigl(2^{p(n)}\bigr)^{B(C_0)} \cdot (\text{combinatorial overhead})
  \;\le\; 2^{p(n) \cdot 2^{2^{p(n)}}}
  \;\le\; 2^{2^{q(n)}}
\]
for some polynomial \(q\).

Each validity clause (V1)--(V10) is a finite-quantifier predicate
over the certificate's components:\@ (V1)--(V3) are local
type/context checks, (V4)--(V5) are graph-reachability checks on
\(S\) (linear in \(|S|\) per query), (V6)--(V8) are mosaic-position
membership and triple checks (polynomial in mosaic count), (V9) is a
congruence check on the explicit-port relation, (V10) is exact
summary computation by reachability.  All in time polynomial in
\(|Q|\), which is polynomial in \(B(C_0)\).
\end{proof}

\begin{theorem}[Effective decidability]\label{thm:decidability}
Concept satisfiability for \(\ALCIRCC\) is decidable by enumerating
finite generated split-forest certificates up to size \(B(C_0) =
2^{2^{p(n)}}\) and checking validity.  Soundness and completeness
are provided by Theorems~6.1 and~7.4 of~\cite{repaired};\@ the
patchwork and quotient steps used inside those theorems are
discharged by Theorems~\ref{thm:patchwork} and~\ref{thm:quotient} of
the present addendum.
\end{theorem}

\begin{proof}
Lemma~\ref{lem:enumeration} enumerates certificates;
Theorem~\ref{thm:patchwork} and Theorem~\ref{thm:quotient} discharge
the only steps in the manuscript's proof of
Theorems~6.1 and~7.4 of~\cite{repaired} that the third-round
review left open;\@ Lemma~7.3 (finite regularisation)
of~\cite{repaired} provides the bound \(|S| \le B(C_0)\) on
representative certificates.
\end{proof}

\begin{remark}[Discharging C8]
With the grammar above and Lemma~\ref{lem:enumeration}, the
obligation \emph{C8 Effective finite certificate language} moves
from \OPEN{} to \CLOSED{}.  Lean target L3 (enumeration bound
mechanization) remains open.
\end{remark}

%% ====================================================================
\section{Status after the addendum}
%% ====================================================================

\begin{center}\small
\renewcommand{\arraystretch}{1.25}
\begin{longtable}{p{0.06\linewidth} p{0.40\linewidth} p{0.16\linewidth} p{0.30\linewidth}}
\toprule
ID & Obligation & Status after addendum & Source \\
\midrule
\endfirsthead
\toprule
ID & Obligation & Status after addendum & Source \\
\midrule
\endhead
\bottomrule
\endfoot
C1 & RCC5 composition table & \PASS{} & WP1 + \cite{repaired}~\S3.1 \\
C2 & Closure and witness thinning & \PARTIAL{} & Hand proof in
\cite{repaired}~Lem.~3.3;\@ Lean-pending \\
C3 & Blocking separated from semantic \(\EQ\) & \PARTIAL{} to \PASS{}
   & (V9) of~\cite{repaired}~Def.~4.5;\@ WP2 corpus \\
C4 & Typed-\(\EQ\) quotient soundness & \CLOSED{} (paper)
   & Theorem~\ref{thm:quotient} of this addendum \\
C5 & Multiple upward \(\PP\)-witnesses & \PASS{}
   & v2~\S3.2 (\(C_{AB}^{\mathrm{inc}}\)) \\
C6 & Request-closed cycles & \PASS{}
   & v2~\S3.3 (pure-isolation test) \\
C7 & Support-closed mosaic patchwork & \CLOSED{} (paper)
   & Theorem~\ref{thm:patchwork} of this addendum + WP6 \\
C8 & Effective finite certificate language & \CLOSED{} (paper)
   & Lemma~\ref{lem:enumeration} of this addendum \\
\end{longtable}
\end{center}

\paragraph{What remains open.}  At paper level only C2 is still
\PARTIAL{}, and only because the witness-thinning lemma of
\cite{repaired}~Lem.~3.3 is hand-proved and would benefit from a
fully expanded proof.  All other obligations are either \PASS{}
(computational) or \CLOSED{} (paper).  The Lean targets L1--L7
from the verification recommendations remain open, in the
priority order L1, L3, L4, L6, then L2/L5/L7.

\paragraph{Honest assessment.}  The addendum upgrades C4/C7/C8 from
\emph{computationally supported, theorem pending} to
\emph{computationally supported, paper theorem closed, Lean
pending}.  It does not claim Lean-level certainty, and it does not
amplify the manuscript's complexity bound:\@ \(B(C_0) =
2^{2^{p(n)}}\) is unchanged.  The addendum's contribution is
making each load-bearing step explicit so that a Lean
formalisation has a definite target.

%% ====================================================================
\begin{thebibliography}{9}
\bibitem{repaired}
\emph{A Generated Split-Forest Decidability Proof for \(\ALCIRCC\)
Without Automata:\@ A Repaired Containment-Oriented Presentation.}
GPT-5.5, May~2026.  In repository at
\texttt{papers/gpt5.5\_round2/repaired\_split\_forest\_no\_automata\_proof.tex}.

\bibitem{RenzNebel1999}
Jochen Renz and Bernhard Nebel.
\emph{On the Complexity of Qualitative Spatial Reasoning:\@ A Maximal
Tractable Fragment of the Region Connection Calculus.}
Artificial Intelligence 108 (1--2), 1999, pp.~69--123.

\bibitem{wessel2002}
Michael Wessel.
\emph{Qualitative Spatial Reasoning with the ALCI\(_{\mathrm{RCC}}\)
Family of Description Logics.}
Technical report, University of Hamburg, 2002/2003.
\end{thebibliography}

\end{document}
